%!TEX root = ../template.tex
%%%%%%%%%%%%%%%%%%%%%%%%%%%%%%%%%%%%%%%%%%%%%%%%%%%%%%%%%%%%%%%%%%%%
%% chapter3.tex
%% NOVA thesis document file
%%
%% Methodology
%%%%%%%%%%%%%%%%%%%%%%%%%%%%%%%%%%%%%%%%%%%%%%%%%%%%%%%%%%%%%%%%%%%%

\typeout{NT FILE chapter3.tex}%

\chapter{Methodology}
\label{cha:methodology}

This chapter describes how the research question of Section~\ref{sec:intro_problem} was
investigated. That question asks whether the apparent robustness of MADDPG routing
controllers to FGSM survives a theoretically stronger observation-space attack, one that
coordinates its perturbation across agents and divides it strategically across time. It
is broken down
into three questions: whether dividing the attack across time buys more damage per
attacked step (RQ1), whether coordinating the perturbation across agents increases the
damage (RQ2), and which architectural choices of the victim determine its vulnerability
(RQ3).

The chapter moves from the broadest methodological decisions to the most concrete
procedures. Section~\ref{sec:research_design} sets out the research approach and the
experimental design, and relates each question to the comparison that answers it.
Section~\ref{sec:starting_point} describes the components inherited from earlier work,
and Section~\ref{sec:victim_architectures} the victims under test.
Section~\ref{subsec:pgd_maddpg} defines the threat model and the attack.
Sections~\ref{subsec:eval_protocol} and~\ref{sec:data_analysis} describe how the data
were collected and analysed, and Section~\ref{sec:validity} how the validity and
reliability of the results were protected. Section~\ref{sec:delimitations} states the
delimitations and limitations of the design, and Section~\ref{sec:method_summary}
summarises the methodology and leads into the results.

% ─────────────────────────────────────────────────────────────────────────────
\section{Research Approach and Design}
\label{sec:research_design}

\subsection{Research approach}
\label{subsec:research_approach}

The research questions are comparative and concern magnitudes: how much packet delivery
an attack costs the victim, and whether one attack costs it more than another. The study
is therefore quantitative, and its data are numerical measurements of packet delivery,
routing decisions and attack timing.

It is conducted as a controlled experiment in simulation, for three reasons. First, the
victims exist as policies trained in the simulator of the earlier
work~\cite{amaral2026maddpg-routing}, and the FGSM baseline was measured in the same
simulator, so a like-for-like comparison with that baseline requires the same
environment. Second, a simulator can replay exactly the same traffic under different
attacks. Attacks can then be compared episode by episode rather than through averages
over different traffic, which removes the largest source of variance in the comparison
(Section~\ref{subsec:analysis_gap}). Third, corrupting the telemetry of a routing
controller on an operational network could be done neither safely nor repeatably.

\subsection{Experimental design}
\label{subsec:experimental_design}

The study follows a factorial design with a control. Four factors are varied:
\begin{itemize}
  \item \textbf{Victim variant:} the seven MADDPG variants of
        Section~\ref{sec:victim_architectures}.
  \item \textbf{Attack scope:} \emph{independent}, in which each agent's perturbation
        is chosen from that agent's own observation, or \emph{coordinated}, in which the
        perturbations of all compromised agents are chosen together
        (Section~\ref{subsec:coordinated_attack}).
  \item \textbf{Timing budget:} the fraction $b$ of steps on which the attack may act,
        set to $0.10$, $0.15$, $0.20$ or $0.25$, or no budget, in which case the attack
        acts on every step (Section~\ref{subsec:timed_attack}).
  \item \textbf{Domain constraint:} the \emph{FGSM-parity} constraint of the baseline or
        the stricter \emph{full} constraint (Section~\ref{subsec:threat_model_pgd}).
\end{itemize}
One adversary is trained for each combination of variant, scope and domain constraint,
28 in total, and each is evaluated at the four budgets and without a budget, giving 140
evaluated configurations. The budget is varied at evaluation rather than in training, so
that the budgets are compared on the same trained adversary
(Section~\ref{subsec:procedure_training}). Everything else is held fixed across
configurations: the perturbation radius, the set of compromised agents, the topology, the
traffic regime and its load, the absence of link failures, and the traffic seed.

Each configuration is run as three arms on identical traffic: a clean arm without
perturbation, a random control, and the attack. The random control is the control
condition of the experiment. It shares the attack's admissible set and timing gate but
chooses its perturbation at random, so the difference between the two arms is
attributable to the direction the adversary learned. The FGSM baseline, evaluated in the
earlier study in the same simulator and on the same traffic, is the external reference.
Table~\ref{tab:rq_design} relates each research question to the comparison that answers
it, the measure used, and where the result is reported.

\begin{table}[htbp]
  \centering
  \small
  \caption{Research questions and the comparisons that answer them; the last two rows
           are supplementary comparisons reported alongside them. The measures are
           defined in Section~\ref{sec:data_analysis}.}
  \label{tab:rq_design}
  \begin{tabular}{>{\raggedright\arraybackslash}p{1.6cm}
                  >{\raggedright\arraybackslash}p{4.8cm}
                  >{\raggedright\arraybackslash}p{3.7cm}
                  >{\raggedright\arraybackslash}p{2.6cm}}
    \hline
    \textbf{Question} & \textbf{Comparison} & \textbf{Measure} & \textbf{Reported in} \\
    \hline
    Central & Learned adversary against FGSM, under the FGSM-parity constraint, on the
              same episodes
            & $\Delta^{\text{FGSM}}$, Eq.~\eqref{eq:fgsm_comparison}
            & Section~\ref{sec:results_fgsm_comparison} \\
    RQ1     & Each timing budget against no budget, for the same trained adversary
            & $\Delta^{\text{adv}}$ and $\Delta^{\text{adv}}$ per unit of realised
              attack rate
            & Sections~\ref{sec:results_parity_clamp} and~\ref{sec:results_full_clamp} \\
    RQ2     & Coordinated against independent adversary, for the same variant,
              constraint and budget
            & $\Delta^{\text{coord}}$, Eq.~\eqref{eq:coordination_effect}
            & Sections~\ref{subsec:results_parity_coordination_effect}
              and~\ref{subsec:results_coordination_effect} \\
    RQ3     & The seven victim variants against one another
            & $\Delta^{\text{adv}}$ by variant
            & Chapter~\ref{chap:results} \\
    ---     & Routing decisions changed against delivery lost
            & Flip rate, Eq.~\eqref{eq:flip_rate}, against $\Delta^{\text{adv}}$
            & Sections~\ref{subsec:results_parity_flips}
              and~\ref{subsec:results_ind_flips} \\
    ---     & FGSM-parity against full domain constraint
            & $\Delta^{\text{clamp}}$, Eq.~\eqref{eq:clamp_effect}
            & Section~\ref{sec:results_clamp_effect} \\
    \hline
  \end{tabular}
\end{table}

% ─────────────────────────────────────────────────────────────────────────────
\section{Starting Point and Inherited Components}
\label{sec:starting_point}

This work is the third stage of a line of research. Because the research question is
comparative, it depends on the earlier stages remaining unchanged, so this section
states which components are inherited and which are contributed here.

\subsection{The routing system under test}
\label{subsec:base_system}

The victim is a multi-agent routing controller developed in earlier work within the
same research line~\cite{amaral2026maddpg-routing}. It applies Multi-Agent Deep
Deterministic Policy Gradient (MADDPG) to adaptive packet routing under the
paradigm of centralised training with decentralised execution
(CTDE)~\cite{lowe2017multi}: each agent $i$ holds an actor $\pi_{\theta_i}$ that
conditions only on its own local observation, while its critic $Q_{\phi^i}$ may,
during training, observe the joint state and the actions of every other agent.
Conditioning the critic on the joint action is what stabilises learning, since from
any single agent's perspective the environment would otherwise appear
non-stationary as the other agents' policies evolve~\cite{lowe2017multi}. After
training the critics are discarded and each agent acts on local information alone.

The network is a service-provider topology of 86 nodes, 65 switches and 21 traffic
endpoints, joined by 106 links, carrying flow-based traffic. The learning agents are
its $M = 14$ distribution switches, the ingress routers to which the endpoints attach;
the remaining switches forward traffic without a learned policy. Each agent observes a
94-dimensional vector: the available bandwidth to each of up to six neighbours
(zero-padded), its queue depth, the diversity of destinations in its queue and the
normalised time step, a presence flag for each of the $n_d = 21$ destinations, the
mean, maximum and variance of its adjacent link utilisations, and, for every
destination, the bottleneck utilisation of each of its $K = 3$ pre-computed candidate
paths.\footnote{The environment assembles a slightly longer vector and truncates it to
94 entries, which drops the last two path utilisations and a final hops-remaining
feature.} Every feature lies in $[0,1]$. The agent's action scores the $K$ candidate
paths of each destination, and the path with the highest score is taken. The trained
system is thus a distributed controller in which no agent holds a global view, yet the
paths chosen collectively determine end-to-end delivery. That prior work also
investigated critic design and graph-based state encoding, which is the origin of the
variant family of Section~\ref{sec:victim_architectures}.

This thesis does not modify the routing controller. The trained agents are treated
as a fixed system under test: their parameters are frozen, and the attacks interact with
them only through their observations and actions.
\subsection{The FGSM baseline}
\label{subsec:base_fgsm}

The second stage established the adversarial setting. A single-step attack based on
the Fast Gradient Sign Method
(FGSM)~\cite{goodfellow2015explainingharnessingadversarialexamples} was developed
against these agents in a preceding dissertation in the same line of
work~\cite{martins2026fgsm}. It is a gradient-based attack, implemented
independently of the learned adversary of this thesis. Applied to the routing
controller, the victim is each agent's trained actor and the input is that agent's local
observation. The objective weights the actor's score for every candidate path by a
smooth, steeply rising function of that path's bottleneck utilisation, read from the
observation, which is close to one for congested paths and close to zero for uncongested
ones, so that ascending it steers the policy towards congested paths rather than
inducing a misclassification. For the GNN variants the gradient is taken through the encoder, with
the other agents' observations set to zero. The perturbation is applied to every agent
at every time step. The threat model it fixed is \emph{observation-space} rather than
actuation-space: the adversary corrupts what the agents perceive, modelling a
compromised telemetry or monitoring plane, while the network itself continues to evolve
according to its true state.

That stage contributed not only an attack but an evaluation protocol: the stressed
traffic regime, the frozen-victim assumption, a random control, and the metrics by
which an attack is judged. This thesis keeps the traffic regime, the frozen-victim
assumption and the metrics unchanged. It also keeps a random control as an arm, but
draws its perturbation differently, uniformly from the $\varepsilon$-ball rather than at
$\pm\varepsilon$ with a random sign (Section~\ref{subsec:procedure_evaluation}).

\subsection{The learned-adversary scaffold}
\label{subsec:base_learned_adversary}

The attack of this thesis is built on an existing implementation, provided as the
starting point for this project, of the learned observation adversary of the
State-Adversarial MDP of Zhang et al.~\cite{zhang2020robust}
(Section~\ref{subsec:sa_mdp}). It trains a network against the frozen victim to perturb
the agents' observations, learning from the damage its perturbations cause, and plugs
into the same evaluation loop as the FGSM baseline. This work uses it as the vehicle for
its two mechanisms: the SAJA-inspired coordination and the timing described
in Section~\ref{subsec:pgd_maddpg} were built into it, together with the random control
used to score it.

\subsection{Contribution of this thesis}
\label{subsec:contribution_scope}

As Section~\ref{sec:intro_problem} explains, the FGSM baseline's limited effect on
delivery admits two readings: either the controller is genuinely robust, or a
single-step, per-agent adversary is too weak for its failure to count as evidence of
robustness. Separating the two requires an adversary without those two limitations,
compared with the baseline on identical episodes and within the same admissible set.
This thesis contributes that adversary, by adding a coordination mechanism and a timing
mechanism to the implementation of Section~\ref{subsec:base_learned_adversary}
(Section~\ref{subsec:pgd_maddpg}), together with the random control and the evaluation
that score it. Such an adversary is theoretically stronger than FGSM
(Section~\ref{subsec:learned_adversary}), so if it does no more damage than the
baseline, the result strengthens the case that the controller is robust rather than
that the baseline was too weak. How much weight that argument can carry depends on the
size of the effects measured against the variation the experiment cannot control, a
question Section~\ref{subsec:limitations_statistics} returns to.

Because the purpose is comparison rather than a new attack in isolation,
compatibility with the earlier stage is a design constraint rather than a
convenience. The attack operates on the same frozen agents, within the same
$\ell_\infty$ ball, and is scored under the same evaluation protocol as the FGSM
baseline. Any difference in measured damage is then attributable to the attack
itself, rather than to a change in the adversary's capabilities or in the conditions
of measurement. The domain constraint is the one respect in which the evaluation goes
beyond this: the attack is evaluated under the baseline's constraint and additionally
under a stricter one, and only results under the former are directly comparable with
the FGSM study (Section~\ref{subsec:threat_model_pgd}).

% ─────────────────────────────────────────────────────────────────────────────
\section{Victim Architectures Under Test}
\label{sec:victim_architectures}

The attacks are evaluated against the family of trained variants inherited from the
base system~\cite{amaral2026maddpg-routing} rather than against a single model, so
that conclusions about the routing task can be separated from conclusions about one
particular network design, which is what RQ3 requires. The trained weights of every
variant are used as they were delivered by the earlier work, which is also what makes
the comparison with the FGSM study possible. All variants share the same actor: a
multilayer perceptron with hidden layers of 256 and 128 rectified units and a sigmoid
output, mapping the local observation to $n_d$ blocks of $K$ path scores, with the
executed route given by the $\operatorname{arg\,max}$ within each block. They differ along three binary
axes, described below.

% ─────────────────────────────────────────────────────────────────────────────
\subsection{Central (CC) versus local (LC) critic}
\label{subsec:axis_critic_domain}

The two critic domains differ in what the value function is allowed to observe
during training:
\begin{equation}
  \underbrace{Q_{\phi^i}\big(\mathbf{g}, a_1, \dots, a_M\big)}_{\text{CC:\ } |E| + M + M|\mathcal{A}|}
  \qquad\text{versus}\qquad
  \underbrace{Q_{\phi^i}\big(s^{\text{local}}_i, a_i\big)}_{\text{LC:\ } d_{\text{obs}} + |\mathcal{A}|},
  \label{eq:arch_critic_domains}
\end{equation}

\noindent with $d_{\text{obs}} = 94$ the per-agent observation dimension and
$|\mathcal{A}| = n_d K = 63$ the per-agent action dimension. The central critic follows
the canonical MADDPG construction~\cite{lowe2017multi} in conditioning on the joint
action, which removes the non-stationarity each agent would otherwise face as its
peers' policies change. Instead of every agent's observation, it receives a compact
global state $\mathbf{g}$: the available bandwidth on each of the $|E| = 106$ links and
the most common destination in each agent's queue. Its input still grows linearly in the
number of agents. The local critic reduces the scheme to independent actor--critic
learners; it forfeits that guarantee but is far cheaper and its size is independent of
the agent count.

Neither critic is used by the attacks of this thesis: the critics are discarded after
training, and every variant is attacked by the same learned adversary. The distinction
matters because it shapes how the victims' policies were trained, which
Chapter~\ref{chap:results} draws on when interpreting the effect of coordination.

% ─────────────────────────────────────────────────────────────────────────────
\subsection{Simple versus duelling value head}
\label{subsec:axis_value_head}

The simple head passes the concatenated state and action through two rectified
hidden layers and emits a scalar. The duelling head replaces the output layer with
two parallel streams whose outputs are summed, $Q(s,a) = V(h) + A(h)$.

One caveat should be stated. In the original duelling
formulation~\cite{wang2016duelingnetworkarchitecturesdeep} the advantage
stream emits one value \emph{per action}, and subtracting the mean advantage makes
the split identifiable. Here the action is an \emph{input} to the critic rather than
an output index, so both streams emit scalars and no such constraint is applied: any
constant may be shifted between $V$ and $A$ without changing $Q$. The duelling head
is thus an extra parallel projection rather than a value--advantage separation, and
differences it produces should be attributed on that basis.

The consequence matters for Chapter~\ref{chap:results}. Summing two scalar streams gives
the same set of functions as emitting one scalar directly, so the two heads span the
same function class and neither can represent a critic the other cannot. What the extra
parameters change is the trajectory training follows, and therefore the weights the
critic settles on and the actor it trains. Any difference in vulnerability that tracks
the value head is consequently a property of the particular solutions these runs
learned, not of what the head is able to express.

% ─────────────────────────────────────────────────────────────────────────────
\subsection{With and without GNN encoding}
\label{subsec:axis_gnn}

The GNN variants insert a shared encoder between the environment and the actor.
Each agent is a graph node whose features are its full observation. Because the agents
are a strict subset of the network's nodes, the graph is the full physical topology: the
72 remaining nodes, the non-agent switches and the traffic endpoints, are included as
relay nodes with zero features. Two \texttt{GraphConv} layers, with a hidden width of 64,
aggregate neighbour information under a residual connection:
\begin{equation}
  \tilde{\mathbf{X}} = \mathrm{ReLU}\Big(
     \mathrm{GraphConv}_2\big(\mathrm{ReLU}(\mathrm{GraphConv}_1(\mathbf{X}, E)), E\big)
     + \mathbf{X} \Big).
  \label{eq:arch_gnn}
\end{equation}

\noindent The encoder is dimension-preserving, so it is a drop-in pre-processor and
the actor is unchanged between variants, and the residual connection makes it a
near-identity transform at initialisation. The relay nodes act purely as
message-passing intermediaries, so information travels at most two hops through the
topology.

The consequence for the attacks is structural rather than one of degree: without the
encoder an agent's policy reads its own observation alone, whereas with it a
perturbation applied to one agent also reaches the encoded state of any agent within two
hops. In the independent scope the perturbations of nearby agents are therefore not
fully independent in their effect. The learned adversary attacks these variants exactly
as it attacks the others, since it needs no gradient path through the encoder.

% ─────────────────────────────────────────────────────────────────────────────
\subsection{Variant matrix}
\label{subsec:variant_matrix}

\begin{table}[htbp]
  \centering
  \caption[MADDPG victim variants under test]{MADDPG victim variants under test, named
           \texttt{<critic domain>-<value head>[-GNN]}.}
  \label{tab:variant_matrix}
  \begin{tabular}{llll}
    \hline
    \textbf{Variant} & \textbf{Critic} & \textbf{Value head} & \textbf{GNN} \\
    \hline
    CC-Simple        & Central & Simple   & --- \\
    CC-Duelling      & Central & Duelling & --- \\
    CC-Simple-GNN    & Central & Simple   & Yes \\
    CC-Duelling-GNN  & Central & Duelling & Yes \\
    LC-Simple        & Local   & Simple   & --- \\
    LC-Duelling      & Local   & Duelling & --- \\
    LC-Duelling-GNN  & Local   & Duelling & Yes \\
    \hline
  \end{tabular}
\end{table}

Table~\ref{tab:variant_matrix} lists the seven variants. \texttt{CC-Simple} is the
canonical MADDPG formulation with no additional architectural machinery, and each other
variant changes one or two of the three axes relative to it. \texttt{LC-Simple-GNN},
the eighth combination, is omitted, as the local-critic value-head comparison is
already covered by \texttt{LC-Simple}/\texttt{LC-Duelling} and the GNN comparison by
\texttt{LC-Duelling}/\texttt{LC-Duelling-GNN}.

% ─────────────────────────────────────────────────────────────────────────────
\section[Coordinated and Strategically-Timed Attacks on MADDPG Routing Agents][Coordinated and Strategically-Timed Attacks]{Coordinated and Strategically-Timed Attacks on MADDPG Routing Agents}
\label{subsec:pgd_maddpg}

The attack of this thesis is a learned observation adversary in the sense of the
State-Adversarial MDP~\cite{zhang2020robust}, trained against the frozen victim, to
which two mechanisms are added. They go beyond the FGSM baseline in the two respects
identified in Section~\ref{subsec:contribution_scope}:
\begin{itemize}
\item \textbf{(A) Coordinated multi-agent perturbation, inspired by SAJA:} the
perturbations of all compromised agents are produced together, from their joint
observation, and judged by their combined effect on the network, as SAJA judges victim
perturbations through the value of the joint action they
induce~\cite{guo2025sajastateactionjointattack}. Several compromised agents can then be
steered towards the same shared link.
\item \textbf{(B) Strategically-timed, critical-state attack, inspired by PGD and by Lin
et al.:} the attack is divided across time, as PGD divides its budget into smaller
steps~\cite{madry2019deeplearningmodelsresistant}, and is spent at chosen moments, as the
strategically-timed attack of Lin et
al.~\cite{lin2019tacticsadversarialattackdeep} perturbs only the steps at which the
victim's decision matters most. A budget limits how often the adversary may act, and a
critical-state score decides when it does, so that the attack is spent at the moments
when the network's redundancy is thinnest.
\end{itemize}

\noindent The threat model is set out first, since it defines what the attacker is
assumed to know and to control. The two mechanisms and the learned adversary that
realises them follow.

% ─────────────────────────────────────────────────────────────────────────────
\subsection{Threat model and admissible perturbation set}
\label{subsec:threat_model_pgd}

The threat model is based deliberately on the one used for the FGSM baseline, so that a
difference in damage can be attributed to the attack rather than to a change in what the
attacker is allowed to do. It differs in what the attacker needs to know, which is
summarised in Table~\ref{tab:threat_knowledge}, and the attack is additionally evaluated
under a stricter domain constraint.

\subsubsection*{Victim and compromised agents}

The victim policies $\{\pi_{\theta_i}\}$ are \emph{frozen}: their parameters are fixed
throughout, and no victim parameter is ever updated, so the attack measures the
robustness of a fixed trained policy rather than the outcome of an adversarial training
game. Let $\mathcal{C} \subseteq \{1,\dots,M\}$ denote the set of compromised agents and
$N = |\mathcal{C}|$ its size. In every experiment of this thesis all agents are
compromised, so $\mathcal{C}$ contains the $M = 14$ distribution switches and $N = 14$.

\subsubsection*{Capability}

The attacker controls the observations of the compromised agents and nothing else.
Whenever it acts, it replaces each compromised agent's observation with a perturbed one,
modelling a compromised telemetry or monitoring plane, and the agents choose their
actions from the perturbed observations. It cannot alter the agents' actions, their
parameters, the traffic, the links or the rewards, and the network continues to evolve
from its true state. When the timing mechanism of Section~\ref{subsec:timed_attack} is
used, the attacker acts on all compromised agents at the same steps and leaves the
observations untouched at the others. For agent $i$ at time $t$ the admissible set is
\begin{equation}
  \mathcal{S}\big(s_{i,t}, \varepsilon\big) =
  \Big\{ \hat{s} \;\Big|\;
    \big\| \hat{s} - s_{i,t} \big\|_{\infty} \le \varepsilon,
    \;\; \hat{s} \in \mathcal{D}
  \Big\},
  \label{eq:admissible_set}
\end{equation}

\noindent with $\varepsilon = 0.30$, where $\mathcal{D}$ is a domain constraint and
$\Pi_{\mathcal{S}}$ denotes the projection onto $\mathcal{S}$: an $\ell_\infty$ clip
about the clean observation followed by a clip of the constrained features onto $[0,1]$.

Two domain constraints are used. The \emph{FGSM-parity} constraint is the one of the FGSM
baseline: it clips the first four features of each observation, which are four of the six
available-bandwidth features, and leaves every other feature free within the
$\varepsilon$-ball. It is needed for the comparison with the baseline, but it allows most
features to leave the range $[0,1]$ in which the environment produces them, so a
perturbed observation can report values that no real telemetry would, such as a negative
link utilisation. The \emph{full} constraint rules this out by clipping every feature onto
$[0,1]$, and so models an attacker whose falsified telemetry stays within the valid
range. It gives a strict subset of the parity set, and its consequences are discussed in
Section~\ref{subsec:learned_adversary}. The random control of
Section~\ref{subsec:procedure_evaluation} always uses the same projection as the attack
it controls, since comparing an attack with a control drawn from a different set would
confound the perturbation direction with the size of the set it is drawn from.

\subsubsection*{Knowledge}

\begin{table}[htbp]
  \centering
  \small
  \caption{What each attack needs to know or access. The learned adversary's
           requirements differ between training and attack time; the timing gate's
           requirement applies only when a timing budget is used.}
  \label{tab:threat_knowledge}
  \begin{tabular}{p{6.4cm}ccc}
    \hline
    & \textbf{FGSM} & \multicolumn{2}{c}{\textbf{Learned adversary}} \\
    \textbf{Knowledge or access} & \textbf{baseline} & \textbf{training} & \textbf{attack time} \\
    \hline
    Victim parameters and gradients                         & Yes & No  & No  \\
    Querying the victim's actions                           & Yes & Yes & No  \\
    Clean observations of the compromised agents            & Yes & Yes & Yes \\
    Write access to those observations                      & Yes & Yes & Yes \\
    A simulator of the network and its traffic              & No  & Yes & No  \\
    Per-step count of dropped and forwarded packets         & No  & Yes & No  \\
    Utilisation of every link (timing gate)                 & No  & Yes & Yes \\
    Available bandwidth of every link and each agent's
    queue destinations (coordinated critic)                 & No  & Yes & No  \\
    \hline
  \end{tabular}
\end{table}

The FGSM baseline is a white-box attack: at every step it differentiates each agent's
actor, and for the GNN variants the encoder, with respect to that agent's observation.
The learned adversary never needs the victim's parameters or gradients. During training
it needs to run the victim inside a faithful simulator of the network, querying the
victim's actions on perturbed observations and receiving, at every step, the number of
packets dropped and forwarded, from which its reward is computed. In the coordinated
scope its critic also reads the available bandwidth of every link and the most common
destination in each agent's queue. At attack time the trained adversary needs only the
clean observations of the compromised agents: in the independent scope each agent's own
observation, and in the coordinated scope all fourteen at once. Whenever a timing budget
is used, however, both in training and at attack time, the timing gate reads the
utilisation of every link in the network to score the current step. The attack is
therefore less demanding than FGSM in what it must know about the victim, but more
demanding in what it must know about the network: it assumes a faithful simulator for
training and, when timed, network-wide link telemetry during the attack.

% ─────────────────────────────────────────────────────────────────────────────
\subsection{(A) Coordinated multi-agent perturbation}
\label{subsec:coordinated_attack}

The coordination mechanism is inspired by SAJA~\cite{guo2025sajastateactionjointattack}
(Section~\ref{subsec:joint_attacks}). SAJA perturbs the observations of several
victim agents at once and scores each candidate perturbation by the centralised
critic's value of the joint action the agents then take, so the perturbations are
judged by their combined effect rather than one agent at a time. This subsection
explains why that matters in a routing network and how the learned adversary implements
it.

\subsubsection*{Separability of the per-agent attack}

The FGSM baseline, and any attack that chooses each agent's perturbation independently,
solves $N$ decoupled problems: for each $i \in \mathcal{C}$,
\begin{equation}
  \eta_{i,t}^{\star} = \operatorname*{arg\,max}_{\eta \,:\, s_{i,t} + \eta \in \mathcal{S}}
  \; \mathcal{L}_{i}\big(s_{i,t} + \eta\big),
  \label{eq:separable_attack}
\end{equation}

\noindent where $\mathcal{L}_i$ depends only on agent $i$'s own observation and
policy. The joint perturbation obtained this way is the concatenation of $N$
independent solutions, and the implied joint objective is \emph{separable},
\begin{equation}
  \mathcal{J}_{\text{sep}}(\hat{\mathbf{s}}_t)
  = \sum_{i \in \mathcal{C}} \mathcal{L}_i(\hat{s}_{i,t}),
  \label{eq:separable_objective}
\end{equation}

\noindent so no term couples two agents. This is a structural limitation rather
than a matter of attack strength: a separable objective cannot express a preference
for outcomes that depend on \emph{which} agents are steered onto the \emph{same}
resource, because the value of steering agent $i$ onto a link is independent of
whether agent $j$ was steered onto it as well. In a topology whose robustness comes
from $K$-path redundancy, the configurations expected to do most damage are those in
which several redirected flows converge on one shared link, and these lie outside the
reach of Equation~\eqref{eq:separable_objective}.

\subsubsection*{Joint perturbation}

The coordinated adversary therefore chooses the perturbations of all compromised agents
together. Writing $\mathbf{s}_t = \big(s_{1,t}, \dots, s_{N,t}\big)$ for the concatenated
clean observations of the compromised agents, a single actor $\nu_\omega$ maps the joint
observation to one joint perturbation direction, and the joint observation presented to
the victim is
\begin{equation}
  \hat{\mathbf{s}}_t = \Pi_{\mathcal{S}}\big( \mathbf{s}_t + \varepsilon\, \nu_\omega(\mathbf{s}_t) \big),
  \qquad \nu_\omega(\mathbf{s}_t) \in [-1,1]^{N d_{\text{obs}}},
  \label{eq:joint_attack}
\end{equation}

\noindent where the projection is applied to each agent's block separately, so that each
agent's perturbed observation lies in its own admissible set. Each agent then receives
its own block of $\hat{\mathbf{s}}_t$, and all compromised agents act on one coordinated
decision rather than on $N$ independent ones. Where SAJA scores a joint perturbation
with the victim's own centralised critic, the coordinated adversary learns its own
critic, which scores the joint perturbation by the damage it is expected to cause and is
additionally given the available bandwidth of every link and each agent's most common
queued destination. Learning its own critic keeps the attack independent of the victim's
internals and makes it applicable to the local-critic variants, which have no joint
critic to use. Because the actor's input and the critic's evaluation are joint, and the
reward is the damage to the whole network, the adversary is not restricted to separable
objectives of the form of Equation~\eqref{eq:separable_objective}.

Algorithm~\ref{alg:coordinated_attack} states the mechanism as it is applied at one time
step. The whole group costs a single actor pass, whose input and output are $N$ times
longer than in the independent scope. Two details matter for what the adversary can
learn. First, the direction stored for learning is the \emph{effective} one, recovered
from the perturbation that survived the projection, so the critic scores what the victim
actually saw rather than the actor's raw output. Second, the step yields one transition
for the whole group, carrying the joint observation, the joint direction and the global
state $\mathbf{g}_t$, whereas the independent scope yields one transition per agent. How
the adversary is trained from those transitions is described in
Section~\ref{subsec:learned_adversary}.

\begin{algorithm}[H]
\SetAlgoLined
\DontPrintSemicolon
\KwIn{clean observations $\{s_{i,t}\}_{i \in \mathcal{C}}$ in group order; actor
      $\nu_\omega$; radius $\varepsilon$; exploration noise $\sigma$ (training only)}
\KwOut{perturbed observations $\{\hat{s}_{i,t}\}_{i \in \mathcal{C}}$; one joint
       transition}
$\mathbf{s}_t \leftarrow \big(s_{1,t}, \dots, s_{N,t}\big)$
  \tcp*{concatenate the group's observations}
$a \leftarrow \nu_\omega(\mathbf{s}_t)$ \tcp*{one actor pass for the whole group}
\lIf{training}{$a \leftarrow \operatorname{clip}\big(a + \xi,\, -1,\, 1\big)$,
  \quad $\xi \sim \mathcal{N}(0, \sigma^2 I)$}
$\hat{\mathbf{s}}_t \leftarrow \Pi_{\mathcal{S}}\big(\mathbf{s}_t + \varepsilon\, a\big)$
  \tcp*{projected per agent block}
$\tilde{a} \leftarrow \big(\hat{\mathbf{s}}_t - \mathbf{s}_t\big) / \varepsilon$
  \tcp*{effective joint direction}
\lForEach{$i \in \mathcal{C}$}{$\hat{s}_{i,t} \leftarrow$ block $i$ of
  $\hat{\mathbf{s}}_t$}
the victim acts on $\{\hat{s}_{i,t}\}$; the network steps from the true state\;
\If{training}{
  $\mathbf{g}_t \leftarrow$ available bandwidth of every link and each agent's most
  common queued destination\;
  store the single transition
  $\big(\mathbf{s}_t,\, \tilde{a},\, r_t,\, \mathbf{s}_{t+1},\, \mathbf{g}_t,\,
  \mathbf{g}_{t+1}\big)$ in $\mathcal{R}$\;
}
\caption{Coordinated joint perturbation at one time step}
\label{alg:coordinated_attack}
\end{algorithm}

% ─────────────────────────────────────────────────────────────────────────────
\subsection{(B) Strategically-timed attacks under an \texorpdfstring{$L_0$}{L0} budget}
\label{subsec:timed_attack}

Without a budget the adversary acts at every time step. A more realistic adversary is
constrained in \emph{how often} it can act, which converts the attack into a
resource-allocation problem over the trajectory.

The mechanism draws on two sources. Dividing the attack in this way is how the design
carries PGD's central idea from the input into time, and spending the resulting budget at
deliberately chosen moments rather than arbitrary ones is the idea of the
strategically-timed attack of Lin et al.~\cite{lin2019tacticsadversarialattackdeep}.
FGSM spends its entire budget in a single step of size
$\varepsilon$, and the FGSM baseline did so at every time step for every agent. PGD
instead divides its budget into a sequence of smaller steps of size
$\alpha$~\cite{madry2019deeplearningmodelsresistant}. By analogy, rather than applying a
full-strength perturbation at every step, the attack divides its effort across the
trajectory and acts on only a fraction of the steps. The analogy concerns \emph{when} the
adversary acts, not how strongly: every step on which it does act is still bounded by
the same $\varepsilon$. The attack is therefore subject to an $L_0$ budget: over a run
of $T$ consecutive decisions and for a budget fraction $b \in (0,1]$,
\begin{equation}
  \frac{1}{T} \sum_{t=1}^{T} \mathbf{1}\big[\text{attack at } t\big]
  \;\le\; b .
  \label{eq:l0_budget}
\end{equation}

\noindent The budget applies to the cumulative attack rate over a whole training or
evaluation run, which spans many episodes, not to each episode separately.

Limiting how often an attack acts, and choosing when to act, was first studied by Lin et
al.~\cite{lin2019tacticsadversarialattackdeep}, whose strategically-timed attack minimises
an agent's expected return subject to a cap of this kind on the number of perturbed
steps, and from whom the name of this mechanism is taken
(Section~\ref{subsec:lit_timed_attacks}). The timing mechanism of this thesis takes that
idea but differs from their attack in three respects. First, its critical-state
score is read from the network, as the utilisation of its links, rather than from the
victim's policy, whose preference between its most and least favoured actions Lin et al.
compare with a threshold. Second, its threshold is recalibrated online as a quantile and
ties at the threshold are rationed, so that the realised attack rate matches the
requested budget, whereas a fixed threshold controls the rate only indirectly. Third, the
perturbation applied at the chosen steps is produced by the learned adversary for all
compromised agents at once, rather than by a Carlini and Wagner attack that drives a
single agent towards its least preferred action.

Spending this budget at random steps would spend much of it at moments when the network
has ample spare capacity. The attack instead spends it at moments of high leverage,
identified by a \emph{critical-state score}, the maximum link utilisation across the
topology,
\begin{equation}
  c_t = \max_{(u,v) \in E} \rho_t(u,v),
  \label{eq:critical_score}
\end{equation}

\noindent where $\rho_t(u,v)$ is the utilisation of link $(u,v)$ at time $t$. This
single quantity responds both to congestion onset, which raises it directly, and to the
loss of a link, which raises it indirectly by concentrating traffic onto the remaining
paths, so the score is elevated exactly when redundancy is thinnest. Computing it
requires the utilisation of every link, which is the network-wide knowledge noted in
Table~\ref{tab:threat_knowledge}.

Because the distribution of $c_t$ is not known in advance, the trigger threshold is
calibrated online. Let $\mathcal{H}_t$ be a sliding window of the $W = 1000$ most recent
scores and
\begin{equation}
  \theta_t = Q_{1-b}\big( \mathcal{H}_t \big),
  \label{eq:timing_threshold}
\end{equation}

\noindent where $Q_{1-b}$ denotes the empirical $(1-b)$-quantile, computed from the
window \emph{before} $c_t$ is inserted so that a step never influences its own decision.
The window is not reset between episodes.

A threshold test of the form $c_t \ge \theta_t$ is not sufficient here, because the
score distribution has an \emph{atom}: a single value taken by a large share of steps.
Link utilisation is accumulated from fixed-size flow reservations, so the score returns
to exactly the same values again and again; the evaluation of
Chapter~\ref{chap:results} locates the atom that matters at the tightest budgets at
$c_t = 0.90$ (Section~\ref{subsec:results_parity_triggers}). Whenever $b$ is smaller
than the fraction of steps scoring at or above an atom, the quantile falls on it,
$\theta_t$ takes the atom's value, and every tied step passes the test. The realised
attack rate is then that fraction rather than $b$: the constraint of
Equation~\eqref{eq:l0_budget} is violated, and any two budgets below that fraction
yield the same attack, so the budget ceases to be the independent variable the
experiment varies. Preliminary runs with the plain threshold test showed exactly this: a
requested budget of $0.10$ realised an attack rate of about $0.21$, and the budgets of
$0.10$ and $0.15$ produced the same attack.

The ties are therefore \emph{rationed}. Writing $D_t$ for the number of gate decisions
up to and including step $t$, and $F_{t-1}$ for the number of attacks already made,
the adversary acts at step $t$ if
\begin{equation}
  \textit{fire}_t =
  \begin{cases}
    1, & c_t > \theta_t, \\[2pt]
    1, & c_t = \theta_t \;\text{ and }\; F_{t-1} + 1 \le b\,D_t, \\[2pt]
    0, & \text{otherwise,}
  \end{cases}
  \label{eq:timing_rule}
\end{equation}

\noindent so that steps strictly above the threshold always fire, while tied steps
are admitted only for as long as the cumulative rate stays within budget. The
tie-break is deterministic rather than randomised: a random rule would introduce
variance between the attacked and the control arm, which are intended to differ only
in the \emph{direction} of the perturbation, and would require a separate random
stream so as not to disturb the seeded traffic sequence.

The first $H = 100$ decisions of a run only observe. They seed the window without
attacking, since the quantile is not meaningful until the window holds enough scores.
Those decisions still count towards $D_t$, so the allowance left unspent during the
warm-up is recovered afterwards. Because steps strictly above the threshold fire
regardless of the running count, the realised rate over a run settles close to $b$ but
can slightly exceed it. Since it is the realised rate, not the requested one, that a
comparison across budgets actually varies, it is measured and reported alongside every
result rather than assumed equal to $b$. The procedure is summarised in
Algorithm~\ref{alg:timing_gate}.

\begin{algorithm}[H]
\SetAlgoLined
\DontPrintSemicolon
\KwIn{budget $b$; window $\mathcal{H}$ of size $W$; warm-up length $H$; running
      counts $D$ (decisions) and $F$ (attacks made); network state at step $t$}
\KwOut{decision $\textit{fire} \in \{\text{true}, \text{false}\}$}
$c_t \leftarrow \max_{(u,v) \in E} \rho_t(u,v)$ \tcp*{critical-state score}
$D \leftarrow D + 1$\;
\eIf{$|\mathcal{H}| < H$}{
  $\textit{fire} \leftarrow \text{false}$ \tcp*{warm-up: observe, do not act}
}{
  $\theta_t \leftarrow Q_{1-b}(\mathcal{H})$ \tcp*{computed before inserting $c_t$}
  \uIf{$c_t > \theta_t$}{
    $\textit{fire} \leftarrow \text{true}$\;
  }\uElseIf{$c_t = \theta_t$}{
    $\textit{fire} \leftarrow (F + 1 \le b\,D)$ \tcp*{ration the tie}
  }\Else{
    $\textit{fire} \leftarrow \text{false}$\;
  }
}
insert $c_t$ into $\mathcal{H}$ (ring buffer of size $W$)\;
\lIf{\textit{fire}}{$F \leftarrow F + 1$}
\Return \textit{fire}\;
\caption{Critical-state timing gate with rationed ties}
\label{alg:timing_gate}
\end{algorithm}

The gate composes with either attack scope: it decides \emph{whether} to act at a
given step, one decision shared by all compromised agents, while the scope decides
\emph{how} the perturbation is produced once the decision to act has been taken. On
skipped steps the agents receive their unmodified observations. When no budget is set
there is no gate, and the adversary acts on every agent at every step.

% ─────────────────────────────────────────────────────────────────────────────
\subsection{Learned observation-space adversary}
\label{subsec:learned_adversary}

Sections~\ref{subsec:coordinated_attack} and~\ref{subsec:timed_attack} describe what
the attack should do. It should perturb the compromised agents jointly, keep each
perturbation within the admissible set, and divide its effort across time by acting
on only a fraction of the steps. Both mechanisms are realised in a learned adversary,
built on the implementation of Section~\ref{subsec:base_learned_adversary}.

The learned adversary comes from the State-Adversarial MDP of Zhang et
al.~\cite{zhang2020robust} (Section~\ref{subsec:sa_mdp}). Because the victim is frozen,
its policies are fixed, which is exactly the condition of Lemma~1: against a fixed
policy, finding the optimal observation adversary is itself an ordinary MDP, and can
therefore be approached with standard reinforcement learning. Such an adversary is
trained here with DDPG, treating the frozen victim as part of its environment. Learning
the perturbation from the damage it causes, rather than solving for it at every step
against an immediate objective, lets the attack value a perturbation by its effect on
later steps and on the network as a whole, and removes the need for the victim's
gradients. The adversary is denoted $\nu_\omega$,
deliberately reusing the symbol of the SA-MDP adversary, with parameters $\omega$.

\subsubsection*{Formulation}

In the independent scope, at time $t$ the adversary observes the clean observation
$s_{i,t}$ of a compromised agent $i \in \mathcal{C}$ and emits a perturbation direction
$a^{\nu}_{i,t} = \nu_\omega(s_{i,t}) \in [-1,1]^{d_{\text{obs}}}$, produced by a $\tanh$
output layer. The observation presented to the victim is
\begin{equation}
  \hat{s}_{i,t} = \Pi_{\mathcal{S}}\big( s_{i,t} + \varepsilon \, a^{\nu}_{i,t} \big),
  \label{eq:la_perturbation}
\end{equation}
so that every perturbation lies in the admissible set by construction. In the
coordinated scope the same construction is applied to the joint observation, as in
Equation~\eqref{eq:joint_attack} and Algorithm~\ref{alg:coordinated_attack}. Lemma~1
takes the perturbed state itself as the
adversary's action; parameterising it instead as a bounded direction scaled by
$\varepsilon$ is an equivalent reparameterisation that keeps the actor's output within
the admissible ball without a penalty term. The victim acts on $\hat{s}_{i,t}$, while the
network continues to evolve from the true state, as the SA-MDP requires.

Lemma~1 prescribes the negative of the victim's reward as the adversary's reward. The
adversary of this thesis departs from this deliberately. The victim is trained on a
shaped per-agent reward combining several utilisation-dependent terms, which is a
training signal rather than the quantity whose degradation this work measures. The
adversary is instead rewarded with the fraction of forwarding events at step $t$ that
end in a drop,
\begin{equation}
  r_t = \frac{n^{\text{drop}}_t}{\max\big(1,\, n^{\text{fwd}}_t\big)},
  \label{eq:la_reward}
\end{equation}
where $n^{\text{fwd}}_t$ counts packets forwarded by any agent at step $t$ and
$n^{\text{drop}}_t$ those dropped. This aligns the adversary's objective with packet
delivery, the outcome measured in Section~\ref{subsec:analysis_measures}. It is a proxy
rather than an exact match: because each packet contributes one forwarding event per
hop, the denominator grows with path length, so the signal is diluted relative to the
end-to-end packet delivery ratio on which the adversary is ultimately scored.

\subsubsection*{Training}

The actor $\nu_\omega$ and a critic $Q^{\nu}_{\psi}$ are multilayer perceptrons with
two hidden layers of 256 ReLU units. The critic scores the perturbation that was
\emph{actually applied}. Because the projection in
Equation~\eqref{eq:la_perturbation} can truncate the actor's output, the stored
direction is the effective one,
$\tilde{a}^{\nu}_{i,t} = (\hat{s}_{i,t} - s_{i,t})/\varepsilon$, rather than the raw
output. Transitions $(s_{i,t}, \tilde{a}^{\nu}_{i,t}, r_t, s_{i,t+1})$ are stored in a
replay buffer $\mathcal{R}$, and the critic is fitted to a one-step temporal-difference
target computed with slowly updated target networks $\nu_{\omega'}$ and
$Q^{\nu}_{\psi'}$,
\begin{equation}
  \mathcal{L}(\psi) = \mathbb{E}_{\mathcal{R}}\Big[
    \big( Q^{\nu}_{\psi}(s, \tilde{a}^{\nu})
      - r - \gamma\, Q^{\nu}_{\psi'}\big(s', \nu_{\omega'}(s')\big) \big)^2 \Big],
  \label{eq:la_critic_loss}
\end{equation}
while the actor ascends the critic's estimate of the damage its perturbation causes,
\begin{equation}
  \max_{\omega} \; \mathbb{E}_{\mathcal{R}}\Big[ Q^{\nu}_{\psi}\big(s, \nu_\omega(s)\big) \Big].
  \label{eq:la_actor_objective}
\end{equation}

In the coordinated scope the same two updates are used, on the joint transitions of
Algorithm~\ref{alg:coordinated_attack} and with the critic additionally conditioned on
the global state: $Q^{\nu}_{\psi}(\mathbf{s}, \tilde{a}, \mathbf{g})$ replaces
$Q^{\nu}_{\psi}(s, \tilde{a}^{\nu})$ in Equations~\eqref{eq:la_critic_loss}
and~\eqref{eq:la_actor_objective}, and the bootstrap term uses
$\mathbf{g}_{t+1}$, the global state after the environment has stepped. Only the critic
sees $\mathbf{g}$: the actor conditions on the joint observation alone, so at attack time
the adversary still needs nothing but the compromised agents' observations
(Table~\ref{tab:threat_knowledge}). This is where coordination enters training. The actor
is trained through a value function that scores a joint perturbation against the state of
the whole network, which is what allows it to prefer perturbations whose damage comes
from several agents being steered onto the same link.

The episode is treated as a continuing task, so the target in
Equation~\eqref{eq:la_critic_loss} is never truncated at the episode boundary. During
training Gaussian noise $\xi \sim \mathcal{N}(0, \sigma^2 I)$ with $\sigma = 0.10$ is
added to the actor's output and clipped back to $[-1,1]$. The remaining settings are
$\gamma = 0.95$, learning rates of $10^{-4}$ for the actor and $10^{-3}$ for the critic,
a minibatch of $m = 256$, a buffer capacity of $2 \times 10^{5}$ transitions, and a
soft target update with rate $\lambda = 5 \times 10^{-3}$. One gradient update is
performed per environment step once $\mathcal{R}$ holds at least
$n_{\text{warm}} = 2000$ transitions. How long each adversary is trained, and under
which conditions, is given in Section~\ref{subsec:procedure_training}. The procedure is
summarised in Algorithm~\ref{alg:learned_adversary}.

\begin{algorithm}[H]
\SetAlgoLined
\DontPrintSemicolon
\KwIn{frozen victim $\{\pi_{\theta_i}\}$; radius $\varepsilon$; timing budget $b$;
      episodes $E$; steps per episode $T_{\text{ep}}$}
\KwOut{adversary actor $\nu_\omega$}
initialise $\nu_\omega$, $Q^{\nu}_{\psi}$; set $\omega' \leftarrow \omega$,
$\psi' \leftarrow \psi$; empty replay buffer $\mathcal{R}$\;
\For{episode $= 1, \dots, E$}{
  reset the network and observe $\{s_{i,1}\}_{i \in \mathcal{C}}$\;
  \For{$t = 1, \dots, T_{\text{ep}}$}{
    \eIf{the timing gate fires at $t$ \emph{(Algorithm~\ref{alg:timing_gate})}}{
      \ForEach{$i \in \mathcal{C}$}{
        $a^{\nu}_{i,t} \leftarrow \operatorname{clip}\big(\nu_\omega(s_{i,t}) + \xi,\, -1,\, 1\big)$,
        \quad $\xi \sim \mathcal{N}(0, \sigma^2 I)$\;
        $\hat{s}_{i,t} \leftarrow \Pi_{\mathcal{S}}\big(s_{i,t} + \varepsilon\, a^{\nu}_{i,t}\big)$;
        \quad $\tilde{a}^{\nu}_{i,t} \leftarrow (\hat{s}_{i,t} - s_{i,t}) / \varepsilon$\;
      }
    }{
      $\hat{s}_{i,t} \leftarrow s_{i,t}$ for every $i \in \mathcal{C}$
      \tcp*{no attack, no transition}
    }
    the victim acts on $\{\hat{s}_{i,t}\}$; the network steps from the true state\;
    $r_t \leftarrow n^{\text{drop}}_t / \max\big(1, n^{\text{fwd}}_t\big)$\;
    \lIf{the gate fired}{store $(s_{i,t}, \tilde{a}^{\nu}_{i,t}, r_t, s_{i,t+1})$
      in $\mathcal{R}$ for each $i \in \mathcal{C}$}
    \If{$|\mathcal{R}| \ge \max(m, n_{\text{warm}})$}{
      sample a minibatch of size $m$ from $\mathcal{R}$\;
      update $\psi$ by Equation~\eqref{eq:la_critic_loss} and $\omega$ by
      Equation~\eqref{eq:la_actor_objective}\;
      $\omega' \leftarrow \lambda\omega + (1-\lambda)\omega'$;
      \quad $\psi' \leftarrow \lambda\psi + (1-\lambda)\psi'$\;
    }
  }
}
\caption{Training the learned observation-space adversary (independent scope)}
\label{alg:learned_adversary}
\end{algorithm}

\subsubsection*{Scope and timing}

Both mechanisms apply to the learned adversary. In the \emph{independent} scope a single
actor, shared by all compromised agents, maps each agent's own observation to that
agent's perturbation. Every agent contributes its own transition, and all transitions
share the network-wide reward $r_t$. Each perturbation is therefore still chosen from one
agent's view alone, as in the separable problem of
Equation~\eqref{eq:separable_attack}, but the shared reward trains the actor towards
perturbations that degrade the network as a whole. In the \emph{coordinated} scope of
Section~\ref{subsec:coordinated_attack}, the per-agent loop of
Algorithm~\ref{alg:learned_adversary} is replaced by
Algorithm~\ref{alg:coordinated_attack}: the observations of all compromised agents are
concatenated into one joint input, a single actor emits one joint perturbation, and the
critic is additionally conditioned on the global state. Everything else in the training
loop is unchanged, including the reward, the timing gate and the update schedule. The
scope therefore changes what one transition contains and how many of them a step
produces: one joint transition for the whole group, rather than one per agent.

The timing gate of Section~\ref{subsec:timed_attack} is applied during training as well
as evaluation. On a step it skips, the victim receives clean observations and no
transition is stored, but the gradient update of Algorithm~\ref{alg:learned_adversary}
still runs. A tighter budget therefore reduces the number of fresh transitions without
reducing the number of updates, and each stored transition is replayed correspondingly
more often. The effect is compounded in the coordinated scope, which stores one
transition per step rather than one per agent. At evaluation the actor is used
deterministically, without exploration noise, and it neither queries the victim nor
learns further.

\subsubsection*{Relation to the FGSM baseline}

The learned adversary is theoretically stronger than the FGSM baseline. Under Lemma~1
the optimal SA-MDP adversary is the best of all deterministic perturbation policies
within the admissible set, and FGSM, which computes its perturbation deterministically
from the current observation, is one such policy. Within the same admissible set and
acting on the same steps, an optimal learned adversary can therefore do no worse than
FGSM with respect to its own objective, and it can additionally exploit what FGSM
cannot: perturbations that pay off over later steps, coordination across agents, and
timing. A trained adversary only approximates that optimum, since it is limited by its
training, which is why the comparison of Section~\ref{sec:results_fgsm_comparison} is
needed. If the theoretically stronger attacker does no more damage than FGSM, the
result supports the robustness of the controller rather than the weakness of the
baseline, provided the damage both attacks do is measured well enough to tell them
apart (Section~\ref{subsec:limitations_statistics}).

It also requires strictly \emph{less} knowledge of the victim, as set out in
Section~\ref{subsec:threat_model_pgd}: FGSM differentiates the victim with respect to its
input observations, whereas the learned adversary never computes a gradient of the
victim. In exchange it needs a faithful simulator of the network during training and,
when timed, network-wide link utilisation during the attack.

Its admissible set depends on the domain constraint. Under the \emph{FGSM-parity}
constraint it perturbs within exactly the set of the FGSM baseline. Under the \emph{full}
constraint every feature is clipped onto $[0,1]$, which yields a strict subset of that
set. Measured on sampled on-policy observations, the full constraint removes roughly a
third of the available perturbation range and binds on around four fifths of
feature--step pairs. Nearly all of this loss comes from the lower limit of $[0,1]$: over
half of the observation vector sits at exactly zero, where clipping removes the downward
half of the $\varepsilon$ ball. A separate adversary is trained under each constraint and
always evaluated under the constraint it was trained with, since an actor's direction is
fitted to the projection it was trained under. Only the results under the FGSM-parity
constraint are directly comparable with those of the FGSM study. The smaller set of the
full constraint does not make its results a lower bound on those under the parity
constraint: each has its own trained adversary, and
Section~\ref{sec:results_clamp_effect} shows the two adversaries harming different
victims.

% ─────────────────────────────────────────────────────────────────────────────
\section{Data Collection}
\label{subsec:eval_protocol}

\subsection{Environment and traffic regime}
\label{subsec:procedure_environment}

All training and evaluation runs take place in the simulator of
Section~\ref{subsec:base_system}, in the same stressed traffic regime as the FGSM
baseline: an offered load factor of $2.0$ with a hotspot traffic matrix. Each flow is,
with probability $0.75$, sent from one of eight designated sources (the endpoints
\texttt{BS1}--\texttt{BS4} and \texttt{MECS1}--\texttt{MECS4}) to one of the two
endpoints \texttt{CS1} and \texttt{CS2}, and otherwise between two endpoints chosen
uniformly. The regime is stressed on purpose: at nominal uniform load the network has
ample spare capacity, so routing decisions, and therefore perturbations of them, barely
affect delivery. No link failures are injected. The adversary's training loop does not
inject failures, so evaluating without them keeps every evaluation within the conditions
its adversary was trained under. An episode lasts 256 steps, and the network is reset at
the start of every episode.

\subsection{Training the adversaries}
\label{subsec:procedure_training}

One adversary was trained for each of the 28 combinations of victim variant, attack
scope and domain constraint, with the settings of Section~\ref{subsec:learned_adversary}.
Each was trained for 300 episodes of 256 steps against its frozen victim, with
$\varepsilon = 0.30$, under the domain constraint it would later be evaluated under,
and with the timing gate set to $b = 0.25$. That budget is the loosest in the sweep, so
the tighter evaluation budgets act on a subset of the states encountered during
training; evaluated without a budget, the adversary also acts on states it did not
encounter. Training is not seeded, so repeated training runs would produce different
adversaries, and each configuration was trained once
(Section~\ref{subsec:validity_reliability}).
Each trained adversary is stored with its actor weights, its perturbation radius, its
scope and the domain constraint it was trained under.

\subsection{Evaluation runs}
\label{subsec:procedure_evaluation}

Each trained adversary was evaluated at the four timing budgets and without a budget.
Each of these 140 configurations is evaluated with three arms of 30 episodes each: a
\emph{clean} arm with no perturbation, a \emph{random} control, and the \emph{attack}
under test. The random control draws the perturbation of every feature independently
and uniformly from $[-\varepsilon, \varepsilon]$, applies the same projection as the
attack, and, when a timing budget is set, is gated by its own copy of the same timing
gate. The attack arm uses the trained actor deterministically.

All three arms are run on identical episodes. The traffic sequence is generated from a
common seed, $20240601$, to which the random generators are reset at the start of every
arm, so episode $e$ of one arm and episode $e$ of another differ only in the
perturbation applied. The random control draws from a dedicated random generator so
that it consumes no entropy from the traffic generator, which would otherwise
desynchronise the arms and invalidate the pairing.

The FGSM baseline is not re-run. Its results are taken from the earlier study, which
evaluated it on 15 episodes at the same load and traffic seed, without link failures,
under the FGSM-parity constraint and at $\varepsilon = 0.30$. Because the traffic
sequence is generated from the same seed, the first 15 of the 30 episodes of every
evaluation are the episodes on which FGSM was evaluated. The FGSM study's random control
differs from the one used here: it perturbs every feature by $\pm\varepsilon$ with a
random sign.

\subsection{Recorded data}
\label{subsec:procedure_records}

Each configuration produces one record, which stores:
\begin{itemize}
  \item the configuration itself: victim variant, scope, perturbation radius, timing
        budget, load, number of link failures, traffic seed and domain constraint;
  \item for each arm, the packet delivery ratio of every episode;
  \item for the random control and the attack, the fraction of routing decisions changed
        (Section~\ref{subsec:analysis_measures});
  \item for each gated arm, the timing gate's counts: decisions taken, steps attacked,
        realised attack rate, warm-up decisions, and ties seen, admitted and rejected.
\end{itemize}
The FGSM-parity runs additionally record why the gate fired: for every attacked step,
the case of Equation~\eqref{eq:timing_rule} that admitted it and the congestion band its
score fell in, namely saturated ($c_t \ge 1.00$), near saturation
($0.90 \le c_t < 1.00$), high congestion ($0.75 \le c_t < 0.90$) or moderate
($c_t < 0.75$). This recording was added after the full-constraint runs had been made.

\subsection{Implementation and platform}
\label{subsec:procedure_platform}

The simulator, the victims, the FGSM baseline and the learned adversary are implemented
in Python, with PyTorch for the neural networks, PyTorch Geometric for the GNN encoder
and NetworkX for the topology. The learned adversary exposes the same interface as the
FGSM baseline, so both attacks are scored by the same evaluation loop. All runs were
executed inside a Docker container on a university GPU server, and the tables and
figures of Chapter~\ref{chap:results} are computed directly from the stored records.

% ─────────────────────────────────────────────────────────────────────────────
\section{Data Analysis}
\label{sec:data_analysis}

\subsection{Measures}
\label{subsec:analysis_measures}

The \emph{outcome} of an episode is its end-to-end packet delivery ratio (PDR): the
percentage of the packets injected during the episode that are delivered by its end.
Packets still in transit when the episode ends count as not delivered. The PDR is
computed in the same way in every arm and in the FGSM study.

The \emph{decisions} an attack changes are measured by the action flip rate, the
fraction of per-agent, per-destination $\operatorname{arg\,max}$ path choices that differ between the
clean and the perturbed observations at the same time step,
\begin{equation}
  \mathrm{flip}
  = \frac{1}{|\mathcal{C}|\, n_d\, T_{\text{ep}}}
    \sum_{t}\sum_{i \in \mathcal{C}}\sum_{d}
    \mathbf{1}\Big[
      \operatorname*{arg\,max}_k \pi_{\theta_i}\big(\hat{s}_{i,t}\big)_{d,k}
      \neq
      \operatorname*{arg\,max}_k \pi_{\theta_i}\big(s_{i,t}\big)_{d,k}
    \Big],
  \label{eq:flip_rate}
\end{equation}

\noindent averaged over the episodes of an arm and counted over every step, including
the steps a timing gate skipped, on which no decision changes. The two are reported
separately because they can diverge. In a topology with $K$-path redundancy, a high flip
rate accompanied by an unchanged PDR indicates that the network absorbed the redirected
flows rather than that the attack succeeded, and reporting a flip rate as though it were
a measure of damage would conflate the two. How far they diverge in practice is reported
alongside the gaps in Chapter~\ref{chap:results}.

The \emph{timing} of an attack is described by its realised attack rate, the fraction of
gate decisions on which it attacked, $F_T / D_T$, and by the composition of its triggers,
the share of attacked steps admitted by each case of Equation~\eqref{eq:timing_rule} in
each congestion band.

\subsection{Adversarial gap and confidence intervals}
\label{subsec:analysis_gap}

Any perturbation of the observations disturbs the controller, whether or not its
direction was chosen adversarially, so the drop in PDR relative to the clean arm does
not isolate what the attacker's choice of perturbation contributes. The primary measure
is therefore the \emph{adversarial gap}, the mean paired per-episode difference in PDR
between the random control and the attack,
\begin{equation}
  \Delta^{\text{adv}} = \frac{1}{n} \sum_{e=1}^{n}
    \Big( \mathrm{PDR}^{\text{rand}}_{e} - \mathrm{PDR}^{\text{atk}}_{e} \Big),
  \label{eq:adversarial_gap}
\end{equation}

\noindent in percentage points, with $n = 30$. A positive gap means the attack cost the
victim more delivery than the random control did. It is reported with a paired $95\%$
confidence interval, $\Delta^{\text{adv}} \pm t_{0.975,\,n-1}\, s_d / \sqrt{n}$, where
$s_d$ is the standard deviation of the per-episode differences and
$t_{0.975,29} = 2.045$. Differencing within episode pairs removes the
episode-to-episode variation in offered traffic, which would otherwise dominate the
comparison. The drop from the clean arm to the attack is reported alongside the gap,
but is not used as the measure of attack strength.

A result is treated as evidence of an effect only when its confidence interval excludes
zero, which is equivalent to a two-sided paired $t$-test at the $5\%$ level. An interval
that spans zero is reported as a null result, not as evidence of robustness, since each
configuration can only resolve effects larger than its interval's half-width. The
interval assumes that the per-episode differences are independent and approximately
normally distributed. Every episode starts from a reset network with fresh traffic, but
the timing gate's window and counters carry over from one episode to the next, so the
episodes of a gated arm are not strictly independent. No formal test of normality was
applied; with 30 pairs the $t$-interval is robust to moderate departures from it.

\subsection{Derived comparisons}
\label{subsec:analysis_comparisons}

Three further comparisons are paired in the same way and carry a confidence interval of
the same form.

The \emph{effect of coordination} compares the coordinated and independent adversaries
trained against the same variant, under the same constraint and at the same budget.
Neither the clean arm nor the random control depends on the adversary, so both are the
same in the two scopes, and the difference between the two adversarial gaps reduces to
the per-episode difference between the two attack arms,
\begin{equation}
  \Delta^{\text{coord}}
  = \Delta^{\text{adv}}_{\text{coord}} - \Delta^{\text{adv}}_{\text{ind}}
  = \frac{1}{n} \sum_{e=1}^{n}
    \Big( \mathrm{PDR}^{\text{atk,ind}}_{e} - \mathrm{PDR}^{\text{atk,coord}}_{e} \Big).
  \label{eq:coordination_effect}
\end{equation}
A positive value means coordination increased the damage.

The \emph{effect of the domain constraint} compares the adversaries trained under the
two constraints. Their random controls draw from different sets and therefore differ, so
the comparison is the paired difference between the two adversarial gaps,
\begin{equation}
  \Delta^{\text{clamp}}
  = \Delta^{\text{adv}}_{\text{parity}} - \Delta^{\text{adv}}_{\text{full}}
  = \frac{1}{n} \sum_{e=1}^{n}
    \Big[ \big( \mathrm{PDR}^{\text{rnd,parity}}_{e} - \mathrm{PDR}^{\text{atk,parity}}_{e} \big)
        - \big( \mathrm{PDR}^{\text{rnd,full}}_{e} - \mathrm{PDR}^{\text{atk,full}}_{e} \big) \Big].
  \label{eq:clamp_effect}
\end{equation}
A positive value means the adversary trained under the FGSM-parity constraint did more
damage, relative to its own random control, than the one trained under the full
constraint.

The \emph{comparison with FGSM} uses only the FGSM-parity runs, in which both attacks
perturb within the same admissible set. The two studies used different random controls,
so their adversarial gaps are measured against different references and are not
differenced. The attack arms themselves are compared instead, on the 15 episodes the two
evaluations share,
\begin{equation}
  \Delta^{\text{FGSM}}
  = \frac{1}{15} \sum_{e=1}^{15}
    \Big( \mathrm{PDR}^{\text{atk,FGSM}}_{e} - \mathrm{PDR}^{\text{atk,learned}}_{e} \Big),
  \label{eq:fgsm_comparison}
\end{equation}
with a paired $95\%$ confidence interval using $t_{0.975,14} = 2.145$. A positive value
means the learned adversary did more damage than FGSM. For reference, FGSM's own
adversarial gap against its own random control is reported on the same 15 episodes.

The pairing of all three comparisons rests on the arms concerned running on identical
traffic, which is checked on the records before each comparison is made
(Chapter~\ref{chap:results}).

\subsection{Damage per attacked step}
\label{subsec:analysis_efficiency}

RQ1 asks whether a limited budget buys more damage \emph{per attacked step}. A gated
attack acts on fewer steps than an ungated one, so its total damage is expected to be
lower, and the comparison is made after normalising by the realised attack rate: the
adversarial gap divided by the realised rate gives the damage per unit of attack rate,
and a tighter budget is more efficient when this quantity is larger than when attacking
every step. In the comparison with FGSM, where the adversarial gaps are measured against
different controls, the drop from the clean arm is normalised instead. These normalised
figures are descriptive and carry no confidence interval.

\subsection{Multiple comparisons}
\label{subsec:analysis_multiplicity}

Chapter~\ref{chap:results} reports well over a hundred confidence intervals, and no
correction for multiple comparisons is applied. A few intervals can therefore be
expected to exclude zero by chance, and conclusions are drawn from patterns that hold
across budgets rather than from isolated significant cells
(Section~\ref{subsec:limitations_statistics}).

% ─────────────────────────────────────────────────────────────────────────────
\section{Validity and Reliability}
\label{sec:validity}

\subsection{Internal validity}
\label{subsec:validity_internal}

The design aims to make every difference between arms attributable to the attacker's
choice of perturbation. Several measures serve this aim.
\begin{itemize}
  \item \textbf{A fixed victim.} The victim's parameters are frozen and its gradients
        disabled, and it acts deterministically. The tooling refuses to proceed if a
        variant's trained weights cannot be found, so no adversary can be trained or
        evaluated against an untrained victim.
  \item \textbf{Identical traffic.} The arms share the traffic seed, and the random
        control draws from its own generator. The records confirm the pairing: the
        clean arm's per-episode series is bit-identical across scopes, across domain
        constraints and with the first 15 episodes of the FGSM study, in every
        configuration (Chapter~\ref{chap:results}).
  \item \textbf{A matched control.} The random control uses the same admissible set and
        the same timing gate as the attack. Each arm gates on its own trajectory, and
        the realised rates of both arms are recorded; an evaluation flags any
        configuration in which they differ by more than five percentage points, since
        the arms would then differ in how often they attack and not only in the
        direction of the perturbation. The deterministic tie-break of
        Section~\ref{subsec:timed_attack} adds no randomness between them.
  \item \textbf{A controlled budget.} The realised attack rate is measured rather than
        assumed, and the tie-rationing rule was introduced after preliminary runs had
        shown a plain threshold failing to control it
        (Section~\ref{subsec:timed_attack}).
  \item \textbf{Matched training and evaluation conditions.} Every adversary is stored
        with the scope and domain constraint it was trained under, the evaluation reads
        the scope from it and warns if the constraint differs, and every reported
        evaluation uses the adversary's training constraint.
\end{itemize}

\subsection{Construct validity}
\label{subsec:validity_construct}

Damage is measured as lost delivery relative to a random perturbation from the same set,
not as the raw drop from the clean arm and not as decisions changed, which are reported
separately. One gap remains between what the adversary optimises and what it is scored
on: its reward is the per-hop drop fraction of Equation~\eqref{eq:la_reward}, while it
is scored on end-to-end PDR. A significantly negative adversarial gap is therefore read
as a failure of the trained adversary or of its objective. Similarly, the claim that the
learned adversary is theoretically stronger than FGSM concerns the optimal adversary;
the trained adversaries only approximate it,
which is why the claim is tested empirically in the comparison with FGSM.

\subsection{Reliability}
\label{subsec:validity_reliability}

Given a trained adversary and the traffic seed, an evaluation is designed to be
reproducible: the traffic is seeded, the victim and the attack's actor act
deterministically, and the random control draws from a seeded generator. The clean arms,
bit-identical across separate evaluations, are consistent with this. Training is not
reproducible in the same way. The adversary's weights are initialised
at random, exploration noise is added to its perturbations and its updates use randomly
sampled minibatches, and training was not seeded, so repeating the training of a
configuration would produce a different adversary. Each configuration was trained once,
and every result describes that single trained adversary. The confidence intervals of
Section~\ref{subsec:analysis_gap} capture the variation between episodes, not between
training runs. This is the most important limitation of the study, and
Section~\ref{subsec:limitations_statistics} discusses its consequences.

% ─────────────────────────────────────────────────────────────────────────────
\section{Delimitations and Limitations}
\label{sec:delimitations}

The scope of the study was deliberately delimited in the following respects:
\begin{itemize}
  \item Only the observation space is attacked, matching the FGSM baseline's threat
        model; attacks on actions, parameters, traffic or links are out of scope.
  \item A single topology and a single stressed traffic regime are used, without link
        failures, so that the results are comparable with the FGSM study and the
        evaluation stays within the adversaries' training conditions.
  \item A single perturbation radius, $\varepsilon = 0.30$, is used, and all fourteen
        agents are compromised. The budget of the attack is varied in time rather than
        in size, and partial compromise is not studied.
  \item Four timing budgets between $0.10$ and $0.25$ are evaluated, together with the
        ungated condition, and every adversary is trained at $b = 0.25$ only.
  \item Seven of the eight combinations of the victim axes are evaluated.
  \item The coordination and timing mechanisms are realised with a learned adversary.
        Optimisation-based PGD and SAJA attacks are outside the evaluation reported
        here; an iterated PGD attack has since been run in the companion
        study~\cite{martins2026fgsm}, and SAJA remains future work
        (Section~\ref{subsec:future_optimisation}).
\end{itemize}

The study is also subject to limitations it could not control within its design. Each
configuration was trained once and without a seed; all evaluations share one traffic
seed; the comparison with FGSM rests on the 15 episodes on which the earlier study
evaluated it; the reasons why the timing gate fired were recorded only for the
FGSM-parity runs; and the adversary's reward is a proxy for the delivery on which it is
scored. Their consequences for the results are discussed in
Section~\ref{sec:limitations}.

% ─────────────────────────────────────────────────────────────────────────────
\section{Summary}
\label{sec:method_summary}

The research question is answered by a controlled, paired experiment in the simulator
in which the victims were trained and the FGSM baseline was measured. A learned SA-MDP
adversary, extended with SAJA-inspired coordination and timing inspired by PGD and by
Lin et al., is trained
against each frozen victim and scored against a random control drawn from the same
admissible set and gated in the same way, so the damage attributed to it is damage its
learned perturbation caused. Keeping the FGSM study's threat model, traffic and evaluation
protocol makes the comparison with the baseline a like-for-like one, and varying the
budget, the scope, the domain constraint and the victim separates the contribution of
each mechanism and of each architectural choice. Chapter~\ref{chap:results} presents the
results, first under the FGSM-parity constraint, including the comparison with FGSM, and
then under the full constraint.
